\section{Projekt}\label{sec:projects-&-course-projects}
\begin{tabular}{>{\footnotesize\bfseries} p{32pt} >{\footnotesize}p{0.35\textwidth}}
    \href{https://web.archive.org/web/20250613073516/lokate.io}{Lokate} &
    Ett multiplattforms-system för transportplanering riktat till förare och platschefer på byggplatsen.
    Utvecklat med React Native, React (Next.js) och en backend med Spring Boot-mikrotjänster (Kubernetes) som kommunicerar via gRPC och GraphQL.
    \vspace{6px} \\ \\

   % PubF & Systemadministration såsom migrering av tjänster till Google Workspace, konfigurering av DNS och förbättring av bokningssystem byggt i Laravel PHP och Vue.js. \vspace{3px} \\
    \href{https://www.student.chalmers.se/sp/course?course_id=34088}{DAT067 Projekt} &
    En radiostreaming-app byggd med Flutter och Dart-backend för REST-API och WebSocket-streaming av ljud.
    Beställd av LNDN AB för kursen.
    \vspace{6px} \\ \\
    \href{https://www.student.chalmers.se/sp/course?course_id=33536}{DAT356 UX \& RE} &
    Tillämpade UX-design och principer för kravhantering (Requirements Engineering).
    Skapade specifikationer, diagram (Use-case, Context, I*) och itererade på mockups baserat på UI-designprinciper.
    %\vspace{3px} \\
    %\href{https://www.student.chalmers.se/sp/course?course_id=34092}{DAT257 Agile} &
    %Under tiden vi lärde oss den agila metoden SCRUM skapade vi en React TypeScript-webbapp och en enkel Java Spring Boot HTTP-server för att beräkna utsläpp för en resa.
    %\vspace{3px} \\
    %\href{https://www.student.chalmers.se/sp/course?course_id=34004}{TDA384} & Täcker både trådad samtidighet och meddelandeskickning mellan processer där Java respektive Erlang användes. \vspace{3px} \\
    %HD & Fullstack-webbplats med TypeScript Vue.js-webbapp och Kotlin med Ktor HTTP-API som är containeriserade med Docker. \vspace{3px} \\
\end{tabular}
